由 Plancherel 定理和 \(|I|=1\)，从 \(L^2(\mathbb R^n)\) 到 \(L^2(I\times\mathbb R^n)\) 的算子范数上界为 \(\|\varphi\|_\infty\)。
下面证明 \(O(h^{-d/2})\) 量级的界，所考虑情形为 \(p=\infty\)。
取一个极大的 \(c\sqrt h\) 分离网，所在球面为 \(S^{n-1}\)，其中 \(c>0\) 固定且充分小。对半径减半后互不相交的球面球计数，可将中心数控制在 \(O(h^{-d/2})\) 以内。同样的体积比较表明，半径为 \(2c\sqrt h\) 的球冠的重叠次数有统一上界。
在这些球冠上构造光滑单位分解。取支撑于 \(|\omega-\omega_\nu|^2/h\) 的凸起函数，使其在距离不超过 \(c\sqrt h\) 时等于一，再除以它们的和。分母至少为一。将所得截断函数乘以 \(\varphi(\eta)\)。
在坐标 \(\eta=(\sqrt h\,z,\rho)\) 中取 \(\omega_\nu\) 为轴，所得截断函数落在一个固定的紧支撑集上，且 \(\rho\) 有正下界。它关于 \((z,\rho)\) 的任意固定阶导数均一致有界。
事实上，每个相邻中心在此坐标系中的横向坐标和轴向坐标分别为 \(O(\sqrt h)\) 和 \(1+O(h)\)。向量 \(\omega=\eta/|\eta|\) 的横向坐标为 \(\sqrt h\,z/(\rho^2+h|z|^2)^{1/2}\)，将差式有理化后可见 \((\omega\cdot\omega_\nu-1)/h\) 在该紧集上一致光滑。因此，相互重叠的凸起函数的自变量 \(|\omega-\omega_\mu|^2/h\) 的各阶导数一致有界。

从乘子相位中去掉线性部分 \(T\rho/h\)。余下的相位因子由下式控制
\[
 \frac{T}{h}(\sqrt{\rho^2+h|z|^2}-\rho)
 =\frac{T|z|^2}{\sqrt{\rho^2+h|z|^2}+\rho}.
\]
当 \(T\in I\) 时，它在支撑集上的导数一致有界。反复分部积分，得到逆 Fourier 变换在缩放后变量中的 \(L^1\) 范数一致有界。由 Jacobian 换元可知，可逆的线性频率变换保持这一核范数；乘回去掉的线性相位只会平移核。因此各球冠对空间核的贡献为 \(O(1)\)，所用范数为 \(L^1\) 范数。球冠总数给出 \(L^\infty\) 端点界，且关于 \(T\) 一致成立。

在 \(p_c=2n/d\) 处使用 \ref{ms:thm-1-2}，其中 \(N=h^{-1}\)，而 \(\log(2+N)\) 与 \(\log(2/h)\) 可比。对线性算子 \(f\mapsto U(T)P_h f\) 在空间与时空 Lebesgue 空间 \cite{r-32} 之间应用 Riesz--Thorin 插值。当 \(2\le p\le p_c\) 时，临界估计的插值权重为
\[
 \frac{1/2-1/p}{1/2-1/p_c}=n(1-2/p),
\]
因此对数指数为 \(d(1-2/p)\)，且不出现 \(h^{-1}\) 的幂。当 \(p_c\le p\le\infty\) 时，临界估计的权重为 \(p_c/p\)，由此得到对数指数 \(2/p\) 和频率指数 \((d/2)(1-p_c/p)=d/2-n/p\)。这正是所述的指数。
